% This file is part of the Open Source project 'proTirone'
% (c) 2025 Karsten Reincke (https://github.com/protirone/protico.ltx)
% It is distributed under the terms of the creative commons license
% CC-BY-4.0 (= https://creativecommons.org/licenses/by/4.0/)

\begin{frame}[fragile]  
  \frametitle{LF03:09:IPv4-Adressen}
  \vspace{0.2cm}
  \begin{footnotesize}
  \begin{center}
    \renewcommand{\arraystretch}{1.5}
    \begin{tabular}{|c|c||c|c|c|c|}
    \hline
    &  \textcolor{gray}{\tiny{B4}} & \textcolor{gray}{\tiny{B3}} & \textcolor{gray}{\tiny{B2}} 
    & \textcolor{gray}{\tiny{B1}}  \\
	  \hline
    \hline
    IPv4-String & \multicolumn{4}{c|}{"192.168.0.0"} \\
    \hline
	  dezimal aufgelöst & \texttt{192} & \texttt{168} & \texttt{0} & \texttt{0} \\
    \hline
    hexadezimal aufgelöst & \texttt{0xC2} & \texttt{0xA8} & \texttt{0x00} & \texttt{0x00}\\
    \hline
    binär aufgelöst & \texttt{11000000} & \texttt{10101000} & \texttt{00000000} & \texttt{00000000}\\
    \hline
  \end{tabular}
\end{center}
\end{footnotesize}
\end{frame}

\begin{frame}[fragile]  
  \frametitle{LF03:09:IPv4-Adressen:Anzahl}
  \vspace{0.2cm}
  \begin{footnotesize}
  \begin{center}
    \renewcommand{\arraystretch}{1.5}
    \begin{tabular}{|c|c|c|c|c|}
    \hline
    \textcolor{gray}{\tiny{B4}} & \textcolor{gray}{\tiny{B3}} & \textcolor{gray}{\tiny{B2}} 
    & \textcolor{gray}{\tiny{B1}}  \\
	  \hline
    \hline
	  \textcolor{blue}{\texttt{$2^8$}} & \textcolor{blue}{\texttt{$2^8$}} & \textcolor{blue}{\texttt{$2^8$}} & \textcolor{blue}{\texttt{$2^8$}} \\
    \hline
    \multicolumn{4}{|c|}{\textcolor{blue}{\texttt{$2^{(8 + 8 + 8 + 8)}$}}} \\
    \hline
    \multicolumn{4}{|c|}{\textcolor{blue}{\texttt{$= 2^{(32)}$}}} \\
    \hline
    \multicolumn{4}{|c|}{\textcolor{blue}{\texttt{$= 4.294.967.296$}}} \\
    \hline
  \end{tabular}
\end{center}
\end{footnotesize}

\vspace{0.4cm}
\begin{flushright}
\begin{tiny}
ergibt bei $8.000.000.000$ Erdbewohnerinnen\\ $~0.53$ IPv4-Adressen pro Person \texttt{;-)}
\end{tiny}
\end{flushright}

\end{frame}

\begin{frame}[fragile]  
  \frametitle{LF03:09:IPv4-Adressen:Welche ist korrekt?}
  \vspace{1cm}
  \begin{center}
  \begin{minipage}{0.5\textwidth}
  \begin{enumerate}
    \item \texttt{19216800}
    \item \texttt{C2.A8.0.0}
    \item \texttt{192.168.320.0}
  \end{enumerate}
  \end{minipage}
  \end{center}
\end{frame}


\begin{frame}[fragile]  
  \frametitle{LF03:09:IPv4-Netze}
  \vspace{0.2cm}
  \begin{footnotesize}
  \begin{center}
    \renewcommand{\arraystretch}{1.3}
    \begin{tabular}{cr|c|c|c|c|c|}
    \cline{3-6}
    & &  \textcolor{gray}{\tiny{B4}} & \textcolor{gray}{\tiny{B3}} & \textcolor{gray}{\tiny{B2}} 
    &  \textcolor{gray}{\tiny{B1}}  \\
	  \cline{3-6}
	  & IPv4-Adresse & \texttt{192} & \texttt{168} & \texttt{0} & \texttt{42} \\
    \cline{3-6}
    & Subnetzmaske & \texttt{255} & \texttt{255} & \texttt{255} & \texttt{0}\\
    \cline{3-6}\\
    \cline{3-6}
	  & (bin)(IPv4-Adresse) & \textcolor{blue}{\texttt{11000000}} & \textcolor{blue}{\texttt{10101000}} & \textcolor{blue}{\texttt{00000000}} & \textcolor{teal}{\texttt{00101010}} \\
    \& & (bin)(Subnetzmaske) & \textcolor{red}{\texttt{11111111}} & \textcolor{red}{\texttt{11111111}} & \textcolor{red}{\texttt{11111111}} & \texttt{00000000}\\
    \cline{3-6}
    = & (bin)(Netzadresse) & \textcolor{violet}{\texttt{11000000}} & \textcolor{violet}{\texttt{10101000}} & \textcolor{violet}{\texttt{00000000}} & \texttt{00000000} \\
   
    & Netzadresse & \textcolor{violet}{\texttt{192}} & \textcolor{violet}{\texttt{168}} & \textcolor{violet}{\texttt{0}} & \texttt{0} \\
    \cline{3-6}\\
    \cline{3-6}
    & Broadcastadresse & \texttt{192} & \texttt{168} & \texttt{0} & \texttt{255} \\
    \cline{3-6}
    & Gatewayadresse & \texttt{192} & \texttt{168} & \texttt{0} & \texttt{1} \\
    \cline{3-6}

  \end{tabular}
\end{center}
\end{footnotesize}
\end{frame}

\begin{frame}[fragile]  
  \frametitle{LF03:09:Exkurs:Notationsumwandlung}
  \begin{footnotesize}
  \begin{columns}
    \column{0.3\textwidth}
    \textbf{dec $\Rightarrow$ bin}
  
    \vspace{0.2cm}
    \texttt{\textcolor{blue}{\tiny{0.}} \textcolor{teal}{43} : 2 = 21 \% \textcolor{red}{1}} $\uparrow$\\
    \texttt{\textcolor{blue}{\tiny{1.}} 21 : 2 = 10 \% \textcolor{red}{1}} $\uparrow$\\
    \texttt{\textcolor{blue}{\tiny{2.}} 10 : 2 = \ 5 \% \textcolor{red}{0}} $\uparrow$\\
    \texttt{\textcolor{blue}{\tiny{3.}} \ 5 : 2 = \ 2 \% \textcolor{red}{1}} $\uparrow$\\
    \texttt{\textcolor{blue}{\tiny{4.}} \ 2 : 2 = \ 1 \% \textcolor{red}{0}} $\uparrow$\\
    \texttt{\textcolor{blue}{\tiny{5.}} \ 1 : 2 = \ 0 \% \textcolor{red}{1}} $\uparrow$ \\

    \column{0.7\textwidth}
    
    \renewcommand{\arraystretch}{1.5}
    \begin{tabular}{c c c c c c c c c}
    \textit{\textcolor{blue}{7.}} &\textit{\textcolor{blue}{6.}} &
      \textit{\textcolor{blue}{5.}} & \textit{\textcolor{blue}{4.}} &
        \textit{\textcolor{blue}{3.}} & \textit{\textcolor{blue}{2.}} &
          \textit{\textcolor{blue}{1.}} & \textit{\textcolor{blue}{0.}} & \\
      \texttt{\textcolor{gray}{0}} & \texttt{\textcolor{gray}{0}} & \texttt{\textcolor{teal}{1}} & \texttt{\textcolor{teal}{0}} & 
        \texttt{\textcolor{teal}{1}} & \texttt{\textcolor{teal}{0}} & \texttt{\textcolor{teal}{1}} & \texttt{\textcolor{teal}{1}} & \\
   \end{tabular}
  \end{columns}

  \vspace{0.6cm}
  \rule{0.9\textwidth}{0.5pt}
  
  \vspace{0.6cm}

  \textbf{bin $\Rightarrow$ dec}
  
  \vspace{0.2cm}

  \renewcommand{\arraystretch}{1.5}
    \begin{tabular}{c c c c c c c c c}
    \textit{\textcolor{blue}{7.}} &\textit{\textcolor{blue}{6.}} &
      \textit{\textcolor{blue}{5.}} & \textit{\textcolor{blue}{4.}} &
        \textit{\textcolor{blue}{3.}} & \textit{\textcolor{blue}{2.}} &
          \textit{\textcolor{blue}{1.}} & \textit{\textcolor{blue}{0.}} & \\
      \texttt{0} & \texttt{0} & \texttt{1} & \texttt{0} & \texttt{1} & \texttt{0} & \texttt{1} & \texttt{1} & \\
      \texttt{0*2}$\mathcolor{magenta}{^7}$ & \texttt{0*2}$\mathcolor{magenta}{^6}$ & 
        \texttt{1*2}$\mathcolor{magenta}{^5}$ & \texttt{0*2}$\mathcolor{magenta}{^4}$ & 
         \texttt{1*2}$\mathcolor{magenta}{^3}$ & \texttt{0*2}$\mathcolor{magenta}{^2}$ & 
           \texttt{1*2}$\mathcolor{magenta}{^1}$ & $\texttt{1*2}\mathcolor{magenta}{^0}$ & \\
      & & 32 & + & 8 & + & 2 & + 1 = & \textcolor{teal}{43} \\
   \end{tabular}
  \end{footnotesize}

  \vspace{0.5cm}
  \tiny{\underline{Anregung}: Warum wäre \texttt{45} ein schlechtes Erklärbeispiel, nicht aber die \textit{Weltsuperzahl} \texttt{42}?}
\end{frame}


\begin{frame}[fragile]  
  \frametitle{LF03:09:IPv4-Netze}
  \vspace{0.2cm}
  \begin{footnotesize}
  \begin{center}
    \renewcommand{\arraystretch}{1.3}
    \begin{tabular}{cr|c|c|c|c|c|}
    \cline{3-6}
    & &  \textcolor{gray}{\tiny{B4}} & \textcolor{gray}{\tiny{B3}} & \textcolor{gray}{\tiny{B2}} 
    &  \textcolor{gray}{\tiny{B1}}  \\
	  \cline{3-6}
	  & IPv4-Adresse & \texttt{192} & \texttt{168} & \texttt{0} & \texttt{42} \\
    \cline{3-6}
    & Subnetzmaske & \texttt{255} & \texttt{255} & \texttt{255} & \texttt{0}\\
    \cline{3-6}
  \end{tabular}
\end{center}


\vspace{0.2cm}

  \begin{itemize}
    \item \textit{Subnetzmaske} 'spaltet' \textit{IPv4-Adresse} in 

    \begin{itemize}
      \item \footnotesize{\textbf{Netzanteil} (hier `B4.B3.B2`) und}
      \item \textbf{Hostanteil} (hier `B1`)
    \end{itemize}
    \item $\Rightarrow$
    \begin{itemize}
      \item \textbf{Netzadresse} :- Netzanteil $\circ$ kleinste Zahl im Hostanteil 
      \item \textbf{Broadcastadresse} :-  Netzanteil $\circ$  größte Zahl im Hostanteil 
      \item \textbf{Gatewayadresse}  :- `Netzadresse + 1`  | `Broadcastadresse -1` | ...
    \end{itemize}
  \end{itemize}
\end{footnotesize}
\end{frame}


\begin{frame}[fragile]  
  \frametitle{LF03:09:IPv4-Adressen:CIDR-Notation}

  \vspace{1cm}
\begin{center}
\huge{
  \texttt{\textcolor{blue}{192.168.0.0}/\textcolor{teal}{24}}
}
\end{center}

\end{frame}

\begin{frame}[fragile]  
  \frametitle{LF03:09:IPv4-Adressen:Automatisierung}

    \begin{center}
    \begin{minipage}{0.9\textwidth}
      \inputminted[fontsize=\footnotesize,linenos]{python}{\snpLf/dm-0309-nw-rule-verification.py}
    \end{minipage}
    \end{center}
\end{frame}